\begin{abstract}
AI agents must maintain long-horizon memory over conversational history to provide relevant context for future interactions. We observe that conversational information has natural event and topic based structure that a tree can capture. Once memory is a tree, however, access and updates work differently from their flat counterparts: the retriever must select branches before reading leaves, and the writer must decide where in the hierarchy new facts belong. Both access and updates now involve reasoning through possible trajectories in the tree, and whether the chosen path is correct can only be judged by the final answer to a future question. We propose \textsc{Struct Memory-R1}, a framework that induces tree-structured memory from raw conversation and uses reinforcement learning from question-answering reward to train three components: a \emph{Memory Manager} for tree updates, a \emph{Retrieve Agent} for tree access, and an \emph{Answer Agent} for downstream evaluation. On the LoCoMo benchmark, \textsc{Struct Memory-R1} achieves the highest overall F1 and LLM-judge accuracy among non-oracle methods, with particularly large gains on multi-hop (+8.6 F1) and temporal (+14.1 F1) questions over the flat-memory Memory-R1 baseline.
\end{abstract}
